\chapter{Exophthalmos}

\section*{Definition}
Exophthalmos (also called proptosis) is forward displacement of one or both eyes from the orbit. A clinician measures it and compares it with the person's baseline and population norms; a single millimetre cutoff is not applicable to everyone.

\section*{When to See a Doctor}
See your doctor if:
\begin{itemize}
  \item You notice one or both of your eyes protruding forward
  \item You have double vision or difficulty moving your eyes
  \item You have eye pain, redness of the eyelids, or fever
  \item You experience vision loss or changes in how you see colors
  \item Your eyes become dry, irritated, or do not close fully
\end{itemize}

\section*{How It Develops}
The most common cause is an increase in the amount of tissue inside the bony eye socket. In thyroid eye disease, the immune system attacks the tissues around the eye, causing the muscles and fat behind the eye to swell and push the eyeball forward. This can stretch the optic nerve and affect eye movement.

\section*{Causes}
\begin{itemize}
  \item \textbf{Thyroid eye disease (Graves' orbitopathy):} An autoimmune condition where the body attacks tissues behind the eyes, most often associated with an overactive thyroid (Graves' disease)
  \item \textbf{Orbital cellulitis:} A bacterial infection inside the eye socket, often spreading from the sinuses
  \item \textbf{Orbital tumor:} Growths such as hemangiomas, meningiomas, or optic nerve gliomas
  \item \textbf{Orbital pseudotumor:} Inflammation of the eye socket tissues with no clear cause
  \item \textbf{Vascular abnormalities:} Abnormal connections between arteries and veins behind the eye
  \item \textbf{Injury:} Fracture of the eye socket bones with bleeding or air trapped behind the eye
\end{itemize}

\section*{Related Symptoms}
\begin{itemize}
  \item Eyelids that do not close completely over the eye
  \item A staring appearance
  \item Double vision from restricted eye muscle movement
  \item Swelling around the eyes
  \item Dry, gritty sensation from the cornea being exposed
  \item Eye pain or pressure
\end{itemize}
\textbf{Important:} Thyroid eye disease is the most common cause of eye protrusion in adults. Smoking and uncontrolled thyroid levels can make it significantly worse.

\section*{Medical Evaluation}
\begin{itemize}
  \item Your doctor will measure how far your eyes protrude using an exophthalmometer and track changes over time.
  \item Blood tests will check your thyroid function (TSH, T4, T3) and thyroid antibodies.
  \item CT or MRI scans of the eye sockets may be done to look for tumors, infection, or other causes.
  \item Your vision, color perception, and visual fields will be tested to check for optic nerve compression.
  \item Eye movement tests will check for double vision and restricted muscle function.
\end{itemize}

\section*{Treatment Options}
\begin{itemize}
  \item Thyroid eye disease requires coordinated eye and thyroid care; smoking cessation and control of thyroid dysfunction are important.
  \item Treatment is based on activity and severity and may include lubricants, selenium in selected mild cases, immunomodulatory therapy, or surgery; steroids and radiotherapy are specialist decisions.
  \item Moderate to severe cases may benefit from medications that block the immune attack on the eye tissues.
  \item Orbital cellulitis requires IV antibiotics and sometimes sinus drainage surgery.
  \item Surgery may be needed to decompress the eye socket, lengthen the eyelids, or correct double vision.
  \item Eye muscle surgery is delayed until the condition has been stable for at least 6 months.
\end{itemize}

\section*{Self-Care and Home Management}
\begin{itemize}
  \item Use artificial tears during the day and lubricating eye ointment at night to protect your eyes from dryness.
  \item Wear sunglasses to reduce light sensitivity and protect your eyes from wind and dust.
  \item If you smoke, stop — smoking worsens thyroid eye disease significantly.
  \item Keep your head elevated when sleeping to reduce swelling around the eyes.
  \item Use selenium only if a clinician recommends it, because benefit depends on the setting and excess selenium can be harmful.
  \item Follow up regularly with your eye doctor and endocrinologist to monitor the condition.
\end{itemize}

\section*{References}
\begin{thebibliography}{99}
\bibitem{exophthalmos:ref1} Bahn RS. ``Graves' Ophthalmopathy.'' N Engl J Med. 2010;362(8):726--738.
\bibitem{exophthalmos:ref2} Bartalena L, Kahaly GJ, Baldeschi L, et al. ``The 2021 European Group on Graves' Orbitopathy (EUGOGO) Clinical Practice Guidelines for the Medical Management of Graves' Orbitopathy.'' Eur J Endocrinol. 2021;185(4):G43--G67.
\bibitem{exophthalmos:ref3} Smith TJ, Hegedus L, Douglas RS. ``Role of Insulin-Like Growth Factor-1 Receptor in Graves' Ophthalmopathy.'' N Engl J Med. 2020;382(8):726--738.
\bibitem{exophthalmos:ref4} Weetman AP. ``Graves' Disease.'' N Engl J Med. 2000;343(17):1236--1248.
\bibitem{exophthalmos:ref5} Burch HB, Perros P, Bednarczuk T, et al. ``Management of Thyroid Eye Disease: A Consensus Statement.'' Lancet Diabetes Endocrinol. 2022;10(6):443--456.
\end{thebibliography}
